\section{Introduction}
Multimodal neuroimaging datasets rarely contain every measurement for every visit. MRI characterizes brain structure; FDG, amyloid and tau PET provide physiological or molecular information. ADNI supports longitudinal investigation of these measurements across participants and study phases \citep{weiner2015adni}. Regional summaries support multimodal classification without training image encoders \citep{zhang2011multimodal}.

Class imbalance refers to unequal diagnosis frequencies. Modality imbalance refers to unequal frequencies of available modality combinations, alongside missing inputs. Their intersection creates joint class--modality imbalance: for four modalities, 15 nonempty patterns and three diagnoses define up to 45 cells. A diagnosis that is common overall can still be rare within a particular pattern. Metrics averaged across classes can therefore conceal failures within acquisition patterns. Availability indicators may also encode selection practices. Exposing these indicators makes their role easier to inspect but does not prevent shortcut use \citep{geirhos2020shortcut}.

We ask whether balancing these joint cells helps a neural classifier and makes its errors more complementary to those of a strong gradient-boosted baseline. The proposed availability-aware fusion (AVF) feeds the same standardized regional blocks, feature-missing indicators and availability bits to LightGBM and a masked neural network. The neural branch uses inverse-square-root class--pattern weights and modality dropout. A scalar mixture combines the two branches' seed-averaged predictions without reconstructing unavailable scans.

The study addresses three questions. First, how do class weighting, joint weighting and modality dropout change discrimination and recall? Second, does fusion improve diagnosis-level and cell-level metrics over seed-ensembled LightGBM? Third, how much of the gain also appears when other neural members are fused with the tree? The complete factorial and member controls are reported alongside the proposed configuration.

The contribution is a task-specific experimental analysis and fusion design built from established components. The principal evidence concerns macro-F1 and supported-cell recall in a grouped development evaluation. The small AUROC change, nonnested selection and absence of an independent final test constrain the conclusions. We do not claim a new attention operator or superiority over official image-based systems.
